\documentclass[11pt]{article}
\usepackage[margin=1in]{geometry}
\usepackage[T1]{fontenc}
\usepackage{lmodern}
\usepackage{amsmath,amssymb,amsthm,mathtools}
\usepackage{graphicx}
\usepackage{hyperref}
\usepackage{microtype}
\usepackage{enumitem}

\newtheorem{theorem}{Theorem}
\newtheorem{lemma}{Lemma}
\newcommand{\R}{\mathbb R}
\newcommand{\C}{\mathcal C}
\newcommand{\ew}{\mathrm{ew}}
\newcommand{\rank}{\operatorname{rank}}
\newcommand{\vecop}{\operatorname{vec}}

\title{A complete $n=3$ answer for the $K_nPK_n$ extreme-ray problem}
\author{Solution packet for arXiv:2006.14543}
\date{21 July 2026}

\begin{document}
\maketitle

\noindent\textbf{Status.} \texttt{partial\_solution\_likely\_valid}.
This packet completely settles the first omitted dimension $n=3$ of an
all-dimensions open problem in Mueller-Hermes~\cite{MH}.  The proof is an exact
finite polyhedral enumeration; it uses no floating-point decisions.  The
general problem for arbitrary $n$ remains open.

\section*{Source question}

The question occurs in Section~9, ``Conclusion and open questions,'' on
printed page~38 of~\cite{MH}.  Here $E_n$ is the all-ones matrix and
``ew-positive'' means entrywise nonnegative.

\begin{center}
\includegraphics[width=0.96\textwidth]{figures/open_problem_crop.png}
\end{center}

The relevant bullet asks: for
\[
  K_n=\frac{2}{n}E_n-I_n,
\]
which entrywise-nonnegative $P\in M_n(\R)$ also have $K_nPK_n$ entrywise
nonnegative, and how many extreme rays does the resulting polyhedral cone
have?

\section*{Result at a glance}

Put
\[
 K=K_3=\frac23E_3-I_3,
 \qquad
 \C_3=\{P\in M_3(\R):P\geq_{\ew}0,\ KPK\geq_{\ew}0\}.
\]
Let $G$ be the symmetry group generated by independent permutations of rows
and columns, transposition, and the involution $P\mapsto KPK$.  Then $\C_3$
has exactly $78$ extreme rays, arranged into three $G$-orbits.  Representatives
and their images under the involution are
\begin{align*}
P_A&=\begin{pmatrix}0&0&0\\0&1&2\\0&2&4\end{pmatrix},&
KP_AK&=\begin{pmatrix}4&2&0\\2&1&0\\0&0&0\end{pmatrix},\\[3pt]
P_B&=\begin{pmatrix}0&0&0\\0&1&2\\1&0&2\end{pmatrix},&
KP_BK&=\begin{pmatrix}2&2&0\\0&1&0\\1&0&0\end{pmatrix},\\[3pt]
P_C&=\begin{pmatrix}0&0&1\\0&1&0\\1&0&0\end{pmatrix},&
KP_CK&=P_C.
\end{align*}
The corresponding orbit sizes are $36$, $36$, and $6$.

\section*{Proof Intuition}

After vectorizing $3\times3$ matrices, $\C_3$ is a full-dimensional pointed
cone in $\R^9$ cut out by only $18$ rational inequalities: nine coordinates
of $P$ and nine coordinates of $KPK$.  A ray of a nine-dimensional
polyhedral cone is extreme exactly when its active inequalities have rank
eight.  Thus every extreme ray is exposed by some choice of eight among the
$18$ inequalities.

This turns exhaustiveness into a small finite proof.  For each of the
$\binom{18}{8}=43{,}758$ choices, compute the exact rational nullspace.  A
one-dimensional nullspace contributes a candidate precisely when one of its
two orientations satisfies all $18$ inequalities.  Primitive-integer
normalization removes duplicates.  The resulting $78$ vectors all have
active rank eight.  Applying the natural symmetries produces exactly the
three displayed orbits.  What remains open in the source is the uniform
classification for general $n$, not the $n=3$ case established here.

\section*{Formal theorem and proof}

We flatten matrices row by row.  Since $K=K^T$, this convention gives
\[
 \vecop(KPK)=(K\otimes K)\vecop(P).
\]
Set
\[
 T=K\otimes K,\qquad
 A=\begin{pmatrix}I_9\\T\end{pmatrix}\in M_{18\times9}(\R).
\]
Thus $\vecop(\C_3)=\{x\in\R^9:Ax\geq0\}$.

\begin{lemma}[Active-rank criterion]\label{lem:rank}
Let $D=\{x\in\R^d:Bx\geq0\}$ be a full-dimensional pointed polyhedral cone.
For nonzero $x\in D$, let $B_0(x)$ consist of the rows $b$ of $B$ with
$bx=0$.  Then $\R_+x$ is an extreme ray of $D$ if and only if
$\rank B_0(x)=d-1$.
\end{lemma}

\begin{proof}
The minimal face of $D$ containing $x$ is obtained by imposing precisely the
equalities active at $x$.  Its linear span is $\ker B_0(x)$, hence its
dimension is $d-\rank B_0(x)$.  The face is a ray exactly when this dimension
is one.  Pointedness excludes a one-dimensional linear subspace.
\end{proof}

\begin{theorem}\label{thm:main}
The cone $\C_3$ has exactly $78$ extreme rays.  Under $G$ they are the three
orbits of $P_A,P_B,P_C$ displayed above, of respective sizes $36,36,6$.
\end{theorem}

\begin{proof}
First, $K^2=I_3$ and $K\mathbf1=\mathbf1$.  The all-ones matrix is therefore
sent to itself and satisfies all $18$ inequalities strictly, so $\C_3$ is
full-dimensional.  It is pointed because $P\geq_{\ew}0$ and
$-P\geq_{\ew}0$ force $P=0$.

We now give an exhaustive finite enumeration.  For every eight-element subset
$I\subset\{1,\ldots,18\}$, form the $8\times9$ submatrix $A_I$.  If
$\ker A_I$ is one-dimensional, choose a nonzero generator $x_I$.  Retain its
ray if either $Ax_I\geq0$ or $A(-x_I)\geq0$, and normalize the retained
orientation to a primitive integer vector.  Deduplication gives $78$ rays.

This procedure is exhaustive: by Lemma~\ref{lem:rank}, every extreme ray has
eight linearly independent active rows, hence occurs for at least one of the
subsets $I$.  Conversely, every retained nonzero vector has at least the eight
independent active rows in $I$; its full active matrix cannot have rank nine,
because it annihilates that nonzero vector.  Its active rank is therefore
eight, so Lemma~\ref{lem:rank} shows that every retained ray is extreme.

All arithmetic in this enumeration is rational.  The verifier
\texttt{code/verify\_n3\_extreme\_rays.py} implements exactly the preceding
loop using SymPy nullspaces.  It tests all $43{,}758$ subsets, reports $78$
primitive rays, and then generates $G$ using row and column permutations,
transposition, and $T$.  The rays split into disjoint orbits represented by
$P_A,P_B,P_C$, with sizes $36,36,6$.  Their ranks are respectively $1,2,3$,
so the three orbits are distinct; rank is preserved by every generator of
$G$.  Direct exact multiplication gives the three displayed matrices $KP_iK$.
The active-constraint ranks reported for all three representatives are eight.
Since $36+36+6=78$, these orbits exhaust the list.
\end{proof}

\section*{Verification report}

Run, from the packet directory,
\begin{verbatim}
conda run --no-capture-output -n sandbox python \
  code/verify_n3_extreme_rays.py
\end{verbatim}
The expected summary is
\begin{verbatim}
active_bases_tested=43758
extreme_rays=78
symmetry_orbits=3
orbit sizes: 36, 36, 6
active ranks: 8, 8, 8
\end{verbatim}
The computation constitutes a proof because it exhausts every possible
rank-eight active basis and all nullspace, feasibility, normalization, rank,
and orbit operations are exact over $\mathbb Q$.  It is not a sampling or
floating-point polyhedral calculation.  The short script is included so that
each finite step can be independently audited.

As hand checks, $K^2=I$, each displayed $P_i$ and $KP_iK$ is entrywise
nonnegative, and the three support-size pairs are $(4,4),(4,4),(3,3)$.
The verifier reports respectively $10,10,12$ active inequalities and active
rank eight in every case.

\section*{Scope, novelty, and human review}

\paragraph{Scope.}
The theorem completely answers the source's matrix-cone question only for
$n=3$.  It does not classify arbitrary $n$, the Pauli PPT cone for $N\geq3$,
or the separate tensor-power questions on the same source page.

\paragraph{Novelty check.}
The run indexes had no result for arXiv:2006.14543 or this cone.  A bounded
current search used the exact source wording, the expressions $K_nPK_n$ and
``Pauli PPT spectra,'' the paper title and author, and later arXiv papers and
citations concerning its extremal-ray problem.  It found the source and
general later work on positive-map tensor powers, but no $n=3$ classification
or general answer to this exact matrix-cone question.  Novelty confidence is
moderate because the search is bounded and the result is an elementary finite
polyhedral classification.

\paragraph{Human-review recommendation.}
Review as an exact computer-assisted partial solution.  The principal checks
are (i) the row-major vectorization identity giving $T=K\otimes K$,
(ii) the active-rank criterion, and (iii) that the verifier truly iterates all
$43{,}758$ eight-row subsets using exact arithmetic.  A secondary check is
whether the source authors intended $n$ to range over all positive integers
or only the multi-qubit subsequence; the printed question explicitly says
every $n\in\mathbb N$, so $n=3$ is within its literal scope.

\begin{thebibliography}{1}
\bibitem{MH}
A. Mueller-Hermes,
\emph{Decomposable Pauli diagonal maps and Tensor Squares of Qubit Maps},
arXiv:2006.14543; Journal of Mathematical Physics \textbf{62} (2021), 092204.
\end{thebibliography}

\end{document}
